\appendix
\section{符号与约定}

\begin{center}
\small
\begin{tabular}{@{}ll@{}}
\toprule
符号 & 含义 \\
\midrule
$M$ & 紧致、定向、无边三维流形 \\
$\Sigma$ & 闭定向曲面（$M$ 沿它切开） \\
$G$ & 紧致单 Lie 群，$d=\dim G$，$h^\vee$ 对偶 Coxeter 数（$SU(N)$ 时 $h^\vee=N$） \\
$A$ & $G$-联络，取值于 $\mathfrak g$ 的一形式 \\
$F$ & 曲率，$F=\dd A+A\wedge A$ \\
$k$ & Chern--Simons 耦合常数（整数） \\
$L$ & Chern--Simons 作用量 (2.1) \\
$W_R(C)$ & 沿 $C$、表示 $R$ 的 Wilson 线 \\
$\cM$ & $\Sigma$ 上平坦联络模空间 \\
$\cM_J$ & 复结构 $J$ 下的 Kähler 模空间 \\
$\cL$ & 行列式线丛（预量子化线丛） \\
$\cH_\Sigma$ & 沿 $\Sigma$ 切开的量子 Hilbert 空间 \\
$H^0(\cM_J,\cL^{\otimes k})$ & Kähler 极化下的量子态空间 \\
$R_\lambda$, $\lambda$ & 最高权为 $\lambda$ 的不可约表示 \\
$S_{ij}$ & 模变换矩阵（Verlinde 公式中的 $S$ 矩阵） \\
$N_{ijk}$ & 融合系数（三点函数维数） \\
\bottomrule
\end{tabular}
\end{center}

\noindent\textbf{三个容易混的约定。}
\begin{enumerate}[leftmargin=*]
  \item 本讲的 $L$ 与河野俊丈书里的 $\mathrm{CS}$ 差一个常数因子：
        $L=2\pi k\,\mathrm{CS}_{\text{河野}}$。比较 $k$ 之前先对齐归一化。
  \item 半形式修正把 $k$ 平移为 $k+h^\vee$。本讲凡写 $q=e^{2\pi\ii/(k+h^\vee)}$ 处
        都默认已含此修正；若只用到 $\cL^{\otimes k}$，则 $k$ 是“裸”的。
  \item 边注 \Wp{370} 指 Witten 论文的 CMP 印刷页；中译本的 PDF 页是重排的，
        不能直接换算，请用页边注“原文 p.N”定位。
\end{enumerate}

\section{练习提示}

\begin{enumerate}[leftmargin=*]
  \item \textbf{三叶结。} 对三叶结的标准投影连续用三次拆接关系，最后归到平凡纽结；
        注意每次光滑化后出现的链环分量数。
  \item \textbf{平凡链环。} 对 $n$ 个分量的平凡链环反复使用拆接关系，得到递推式
        $J_n=(-t^{1/2}-t^{-1/2})J_{n-1}$。
  \item \textbf{大规范变换。} 把 $g$ 写成 $S^3$ 上的映射，用度数公式
        \[
        \deg g=\frac{1}{24\pi^2}\int\Tr\left((g^{-1}\dd g)^3\right).
        \]
  \item \textbf{泊松括号。} 把 (3.1) 写成 $\int p^a\dot q^a$ 的形式，
        其中 $p,q$ 是 $A$ 的两个分量；正则括号即得 (3.2)。
  \item \textbf{截面计数。} 用 Riemann--Roch：$\dim H^0=\deg L+1-g$（$g=1$ 时 $\deg L=k$）；
        半形式把 $\cL^{\otimes k}$ 换成 $\cL^{\otimes k}\otimes K^{1/2}$，度数加 $h^\vee$ 的相应倍数。
  \item \textbf{$S$ 矩阵。} 对 $SU(2)$、$k=1$：$S=\frac{1}{\sqrt2}\begin{pmatrix}1&1\\1&-1\end{pmatrix}$，
        $S^2=1$ 直接验证。
  \item \textbf{$Z(S^3)$。} 它等于
        \[
        S_{00}=\sqrt{\frac{2}{k+2}}\,\sin\frac{\pi}{k+2}.
        \]
        取 $k=1,2,3$ 时分别为 $1/\sqrt2$、$1/2$、$\sqrt{2/5}\sin(\pi/5)$。
  \item \textbf{融合系数。} 用 (7.4) 取 $i=j=k=1/2$，$m=0$，再用 $S$ 的显式值。
\end{enumerate}

\section{原文—讲义对照表}

下表把 Witten 原文的各节与本书的讲次对应起来，方便读者两边对照。
页码是 CMP 印刷页；“公式”列给出该节的关键公式编号。

\begin{center}
\small
\begin{tabular}{@{}p{3.4cm}p{1.6cm}p{2.4cm}p{4.4cm}@{}}
\toprule
原文 & 页 & 本书 & 关键内容 \\
\midrule
引言 & 351–353 & 第 1 讲 & Atiyah 两问；Jones 缺三维定义 \\
\S1 作用量 & 353–356 & 第 2 讲 & (1.3) 作用量；(1.4) $k\in\ZZ$；(1.5)(1.6) \\
\S2 弱耦合 & 356–365 & 第 3、6 讲 & (2.16)(2.17) $\eta$ 修正；$q$ \\
\S2.1 纽结 & 362–365 & 第 5 讲 & 穿刺点与 Wilson 线 \\
\S3 正则量子化 & 366–374 & 第 3、4 讲 & (3.1)–(3.3) 辛形式与约束 \\
\S3.1 全纯观点 & 368–369 & 第 4 讲 & 极化与全纯截面 \\
\S3.2 平坦丛 & 369–371 & 第 4 讲 & 射影平坦联络 \\
\S3.3 Wilson 线 & 371–373 & 第 5 讲 & (3.4) 算符值约束；BWB \\
\S3.4 带标记球面 & 373–374 & 第 6、7 讲 & (3.5)(3.6)；融合系数 \\
\S4 可计算性 & 374–394 & 第 7 讲 & (4.28) 手术；(4.51) Verlinde \\
\S4.5 具体手术 & 388–393 & 第 7 讲 & (4.41) 标架；(4.51) \\
\S4.6 辫群 & 394 & 第 7 讲 & Jones 由迹给出 \\
\S5 物理应用 & 394–397 & 第 6、8 讲 & 共形块配对；两处误读 \\
\bottomrule
\end{tabular}
\end{center}

\section{术语表}

\begin{itemize}[leftmargin=*]
  \item \textbf{预量子化（prequantization）}：由整性条件 $[\omega]/2\pi\in H^2(M,\ZZ)$ 决定的线丛与联络。
  \item \textbf{极化（polarization）}：复化切丛的一个可积拉格朗日子丛；用来把 $L^2$ 截面砍成全纯截面。
  \item \textbf{半形式修正（metaplectic correction）}：额外乘 $K^{1/2}$，使能级平移 $k\to k+h^\vee$。
  \item \textbf{共形块（conformal block）}：带标记点曲面上 WZW 模型的态空间，也是 CS 理论在曲面上的态空间。
  \item \textbf{融合规则（fusion rule）}：两个表示张量积分解成不可约表示的重数 $N_{ijk}$。
  \item \textbf{$S$ 矩阵}：模群生成元 $S$ 在环面共形块上的作用，其矩阵元由 (7.2) 给出。
  \item \textbf{手术（surgery）}：挖去管状邻域再以不同方式粘回，用于构造三维流形。
  \item \textbf{标架反常（framing anomaly）}：配分函数在标架改变下的相位，由 $c$ 与标架整数决定。
  \item \textbf{范畴化（categorification）}：把多项式不变量提升为链复形/同调，Euler 特征退回原不变量。
\end{itemize}

\section{部分练习解答}

下面给出几道题的完整解答，其余留给读者；每道题的“关键一步”都写出来了，
卡住时对照 \S B 的提示。

\subsection*{第 1 讲第 1 题：三叶结}

三叶结的标准投影有三个正交叉。对其中一个交叉用 (1.1)：
\[
t^{-1}J(L_+)=t\,J(L_-)+(t^{1/2}-t^{-1/2})J(L_0).
\]
$L_-$ 是两个分量的平凡链环，其 Jones 多项式为 $-(t^{1/2}+t^{-1/2})$；
$L_0$ 是平凡纽结，值为 $1$。代回并整理，再对剩余交叉重复，得到
\[
J(\text{三叶结})=-t^{-4}+t^{-3}+t^{-1}.
\]
关键在于每一步都要重新判断 $L_-,L_0$ 的类型与分量数，这是最容易出错的地方。

\subsection*{第 3 讲第 4 题：泊松括号}

把 (3.1) 写成 $\int\dd t\int_\Sigma \epsilon^{ij}\Tr(A_i\partial_tA_j)$。
令两个空间分量为 $q^a=A_1^a$、$p^a=(k/4\pi)\epsilon^{ab}A_b$（这里 $a,b$ 是二维指标），
则作用量是 $\int p^a\dot q^a$ 的形式，正则括号 $\{q^a,p^b\}=\delta^{ab}$ 立即给出 (3.2)。
唯一需要注意的是 $\Tr$ 的归一化与 $\epsilon$ 的约定，两者都会影响系数 $4\pi/k$。

\subsection*{第 5 讲第 3 题：截面维数}

$G=SU(2)$、$\Sigma=T^2$、$M_J$ 为枕头套。取 $g=1$，Riemann--Roch 给
$\dim H^0(\cM,\cL^{\otimes k})=\deg(\cL^{\otimes k})=k$。
加上半形式 $K^{1/2}=\cL^{-2}$，度数减去 $2$ 的相应倍数后得 $k+1$。
这与“有效 level $k+2$ 的可积表示数减一”一致：
$k+2$ 个表示 $j=0,\frac12,\dots,\frac{k+1}{2}$，减去真空给出的那个平凡维数。

\subsection*{第 7 讲第 2 题：$Z(S^3)$}

$S^3$ 可以看作沿平凡纽结做手术（标架 $0$）得到，手术公式给 $Z(S^3)=S_{00}$。
代入 (7.2)：
\[
S_{00}=\sqrt{\frac{2}{k+2}}\sin\frac{\pi(2\cdot0+1)(2\cdot0+1)}{k+2}
=\sqrt{\frac{2}{k+2}}\sin\frac{\pi}{k+2}.
\]
$k=1$：$\sqrt{2/3}\sin(\pi/3)=1/\sqrt2$；$k=2$：$\sqrt{2/4}\sin(\pi/4)=1/2$。
这两个数可以用 (7.4) 与融合系数独立核对。

\subsection*{第 7 讲第 3 题：$k=1$ 的融合系数}

$k=1$ 时可积表示只有 $j=0,\frac12$。由 (7.4) 与 $S$ 的显式值，
$N_{1/2,1/2,0}=1$（$1/2\otimes1/2$ 含真空），$N_{1/2,1/2,1/2}=0$（$j=1$ 不可积）。
用 $S$ 的公式验证：
\[
N_{1/2,1/2,0}=\sum_{l=0,1/2}\frac{S_{1/2,l}S_{1/2,l}S_{0,l}}{S_{0,l}}
=\sum_l \frac{S_{1/2,l}^2}{1}=1,
\]
因为 $S_{1/2,0}^2=S_{1/2,1/2}^2=\frac12$。这正是 Verlinde 公式在最小情形的手算。

\section{复现与工具}

本讲义的源文件是 \texttt{preamble.tex}、\texttt{main.tex}、\texttt{chapters/L0--L8}、
\texttt{chapters/refs.tex} 与 \texttt{chapters/appendix.tex}。编译需要 xelatex
（ctex、fandol、tcolorbox、longtable、enumitem、booktabs、hyperref）：

\begin{quote}
\small\ttfamily
xelatex -interaction=nonstopmode main.tex\\
xelatex -interaction=nonstopmode main.tex
\end{quote}

\noindent 两遍是必要的：第一遍写目录与交叉引用，第二遍才排对。

同一仓库还包含另外三部分：
\begin{itemize}[leftmargin=*]
  \item 注疏版（按原文节序逐句对照）；
  \item 数据校验与锚点检查脚本；
  \item 协作生产流程记录（角色、议程、逐字稿、评分）。
\end{itemize}
\texttt{deliverables/workflow/} 记录本讲义的协作生产流程。

引用原文时，本讲义用 CMP 印刷页（351--399），可在 \S C 的对照表里查到每节位置。
中文译本重排为 40 页，不能直接按 PDF 页换算，请用页边注“原文 p.N”定位。
